% Generated draft: requires mathematical and source review.
\documentclass[11pt]{article}
\usepackage[margin=25mm]{geometry}
\usepackage{fontspec}
\setmainfont{Latin Modern Roman}
\usepackage{amsmath,amssymb,mathtools,unicode-math}
\setmathfont{Latin Modern Math}
\usepackage{enumitem,needspace,xurl,hyperref}
\setlength{\parindent}{0pt}
\setlength{\parskip}{5pt}
\setlength{\emergencystretch}{3em}
\newcommand{\N}{\mathbb{N}}
\newcommand{\Z}{\mathbb{Z}}
\newcommand{\Q}{\mathbb{Q}}
\newcommand{\R}{\mathbb{R}}
\newcommand{\C}{\mathbb{C}}
\newcommand{\F}{\mathbb{F}}
\newcommand{\A}{\mathbb{A}}
\newcommand{\PP}{\mathbb{P}}
\newcommand{\E}{\mathbb{E}}
\newcommand{\eps}{\varepsilon}
\newcommand{\dd}{\,\mathrm{d}}
\begin{document}
\section*{UnsolvedMath ID: 20001867}
\textbf{Source problem code:} catalogue entry
\subsection{Definitions and mathematical statement}
The record concatenates two questions. Problem 18 is attributed to Boumediene Et-Taoui. Let \(v(n,r,\mathbb K)\) be the largest size of a family of distinct equi-isoclinic \(n\)-planes in \(\mathbb K^r\), \(\mathbb K=\R,\C,\mathbb H\). A standard projection formulation is \(P_iP_jP_i=\lambda P_i\) for every \(i\ne j\), with one common \(0\leq\lambda<1\). It asks:
\begin{enumerate}
\item Find \(v(2,2r,\R)\) for every \(r\geq4\).
\item Find \(v(2,7,\R)\).
\item Find \(v(3,3r,\R)\) for every \(r\geq3\).
\item List all regular \(v\)-tuples in the source's \(G(4,4v)\) whose symmetry groups are isomorphic to \(S_v\). The Grassmannian, regularity and effective-symmetry conventions are not supplied.
\item Find \(v(3,3r,\mathbb K)\) for \(\mathbb K=\C,\mathbb H\); the range of \(r\) is not repeated explicitly.
\item Find all complex \((2M,M)\) equiangular tight frames.
\end{enumerate}
Problem 19, attributed to Emily King, asks to minimize coherence among complex unit-norm tight frames when no equiangular tight frame exists. The aid addresses only 18(f) and 19, not the remaining isoclinic classification questions.

For \(N\geq M\geq1\), a complex UNTF is a matrix \(F=[f_1\cdots f_N]\in\C^{M\times N}\) with \(\|f_i\|=1\) and \(FF^*=(N/M)I\). It is an ETF if all \(|\langle f_i,f_j\rangle|\), \(i\ne j\), are equal. Let \(G=F^*F\). The reported exact conference reduction for \(N=2M\) is
\[
 Q=\sqrt{2M-1}(G-I).
\]
Then \(F\) is an ETF if and only if \(Q\) is a Hermitian complex conference matrix: zero diagonal, unit-modulus off-diagonal entries and \(Q^2=(2M-1)I\). Conversely \(G=I+Q/\sqrt{2M-1}\) is a positive semidefinite Gram matrix of rank \(M\). Unitary frame equivalence, column phases and permutations correspond to the associated Gram/conference monomial-conjugacy equivalences. This translates the classification, rather than enumerating all conference matrices.

For Problem 19 define \(\mu(F)=\max_{i<j}|\langle f_i,f_j\rangle|\). The UNTF space is compact, so \(\mu_*\) is attained. In the nontrivial range \(N>M\), \(N\geq3\), set
\[
 \beta=\frac{N-M}{M(N-1)},\qquad K=\binom N2,\qquad
 V(F)=\frac1K\sum_{i<j}(|\langle f_i,f_j\rangle|^2-\beta)^2.
\]
The unordered squared correlations have mean \(\beta\). The aid reports
\[
 \mu(F)^2-\beta\geq
 \max\left\{\frac{V(F)}\beta,
 \sqrt{\frac{V(F)}{K-1}}\right\}.
\]
If no ETF exists, \(\nu=\min_{F\ \mathrm{UNTF}}V(F)>0\), hence
\[
 \mu_*^2\geq\beta+
 \max\left\{\frac\nu\beta,\sqrt{\frac\nu{K-1}}\right\}>\beta.
\]
This is a strict gap above Welch, but a numerical optimum still requires evaluation of \(\nu\) or stronger information.

The aid's row identity is \(\sum_{j\ne i}|\langle f_i,f_j\rangle|^2=(N-1)\beta\). For \(0\leq\tau<\beta\), it yields the active-correlation count
\[
 \#\{j\ne i:|\langle f_i,f_j\rangle|^2>\tau\}
 \geq\left\lceil\frac{(N-1)(\beta-\tau)}{\mu(F)^2-\tau}\right\rceil.
\]
These constraints do not determine the best coherence for every pair \((N,M)\), and the displayed denominators require the stated range rather than degenerate orthonormal cases.
\subsection{Short English statement}
The merged record asks several isoclinic classification questions and a non-ETF coherence optimization. The aid reduces the \((2M,M)\) complex ETF case to Hermitian conference matrices and proves a strict variance-based Welch gap plus rowwise correlation counts for UNTFs. It does not resolve all the concatenated questions.
\subsection{Sources}
\begin{itemize}
\item Local UnsolvedMath registry, permanent ID \(20001867\), code catalogue entry; canonical \texttt{aim-geometry-notes.json}, \texttt{notes[204]}.
\item AIM, ``Open Problems in Frame Theory/Phase Retrieval,'' after the August 2013 workshop ``Frame theory intersects geometry,'' Problems18 and19, attributed to Boumediene Et-Taoui and Emily King, \url{https://aimath.org/pastworkshops/frametheorygeomproblems.pdf}. The supplied verification reports the OCR concatenation and PDF checks; no PDF inspection is claimed here.
\item No additional authoritative reference is supplied for the conference reduction, variance bound or active-set criterion. The verification's reference label [1] is not expanded in this extract.
\item Research attempt 1: ``partial\_result,'' problem status at run ``open,'' registry status ``partially\_solved,'' rigor label ``proved,'' novelty confidence ``low.'' The machine-generated aid is not peer reviewed; independently verify before citation.
\end{itemize}
\end{document}
